\section*{Erd\H{o}s Problem 75: hereditary finite independence and uncountable chromatic number}
\subsection*{1. Statement and scope}
Does a graph on exactly $\aleph_1$ vertices have chromatic number $\aleph_1$ while, for every $\epsilon>0$, all sufficiently large finite subgraphs $H$ satisfy $\alpha(H)>|H|^{1-\epsilon}$? Can the stronger hereditary bound $\alpha(H)\ge c|H|$ hold for some $c>0$? The cardinality restriction is part of the question. We prove an equivalent formulation of the first finite condition and the correct consequence of the linear condition; we do not construct the required uncountable graph.

\subsection*{2. Hereditary independence implies subpolynomial finite colouring}
Suppose that, for some $0<\delta<1$ and integer $N$, every induced subgraph on $t\ge N$ vertices has an independent set of size at least $t^{1-\delta}$. In a finite induced subgraph on $n$ vertices, remove such an independent set as one colour class while at least $N$ vertices remain. If a step removes $d\ge t^{1-\delta}$ of the current $t$ vertices, concavity of $u^\delta$ gives
\[
 t^\delta-(t-d)^\delta\ge\delta t^{\delta-1}d\ge\delta.
\]
The potential starts at $n^\delta$ and stays nonnegative, so there are at most $n^\delta/\delta$ such steps. Colour the fewer than $N$ leftover vertices separately. Hence
\[
 \chi(H)\le n^\delta/\delta+N. \tag{1}
\]
If the condition in the question holds for every exponent, choose $\delta=\min\{\epsilon/2,1/2\}$. Equation (1) is then less than $n^\epsilon$ for all sufficiently large $n$, uniformly over the finite induced subgraph.
Conversely, if for every $\epsilon>0$ all sufficiently large finite induced $H$ have $\chi(H)\le |H|^\epsilon$, choose a largest colour class with exponent $\epsilon/2$. It has size at least $n^{1-\epsilon/2}>n^{1-\epsilon}$ for $n>1$. Thus the first finite condition is equivalent to a uniform subpolynomial bound on finite induced chromatic numbers. The word hereditary is indispensable: (1) uses the hypothesis again after every removal.

\subsection*{3. The linear variant}
If every sufficiently large induced subgraph has an independent set of size at least $ct$ with $0<c<1$, the same procedure leaves at most $(1-c)^j n$ vertices after $j$ steps. Consequently
\[
 \chi(H)\le\frac{\log n}{-\log(1-c)}+O_{c,N}(1).
\]
For $c=1$, the finite condition forces the graph to be edgeless: extend any alleged edge to a sufficiently large finite vertex set. These bounds are not a fixed finite bound on all finite chromatic numbers, so the finite-colour compactness theorem does not turn them into a finite colouring of the entire graph. They also supply no countable colouring.

\subsection*{4. Correct directions of the elementary inequalities}
For a single graph, $\chi(H)\ge n/\alpha(H)$ is a lower bound on $\chi$, and a large independence number alone gives no upper bound on $\chi$. For example, a clique on $\lfloor n/2\rfloor$ vertices together with the other vertices isolated has $\alpha(H)\ge n/2$ and $\chi(H)=\lfloor n/2\rfloor$. The large independent set cannot be used repeatedly after it has been removed.
Likewise the Caro--Wei estimate $\alpha(H)\ge n^2/(2e(H)+n)$ is sufficient for finding independent sets in sparse graphs; it cannot be reversed to deduce sparsity from large $\alpha$. A complete bipartite graph has $\alpha(H)\ge n/2$ and may have $n^2/4$ edges. Indeed every finite induced subgraph of an infinite complete bipartite graph has an independent set of at least half its vertices. This meets the finite independence condition with dense finite subgraphs, although its chromatic number is only two. Both reversed implications in the earlier imported attempt are therefore invalid.

\subsection*{5. Assessment and attribution}
The finite equivalence and the logarithmic consequence for the linear variant are rigorous reductions. What remains unresolved is the existence of a graph with these hereditary finite properties, exactly $\aleph_1$ vertices and no countable proper colouring. Finite sparsity is one possible sufficient strategy, not a necessary condition proved here. This replaces an erroneous GPT-5.2 bundle with corrected arguments; no solution or new literature result is claimed.
Source: \url{https://www.erdosproblems.com/75}.

\textbf{Completion rate estimate: 15\%.}
